\begin{table}[!ht]
\centering
\caption{H1 --- Inflation broadens into the high tail}
\label{tab:h1}
\setlength{\tabcolsep}{6pt}\footnotesize
\begin{threeparttable}
\begin{tabular}{c c c c}
\multicolumn{4}{l}{\textit{Panel A. Monthly total connectedness by quantile}}\\
\toprule
$\tau$ & Overall & Contemporaneous & Lagged \\
\midrule
0.1 & 5.5 & 2.8 & 2.8 \\
0.3 & 7.5 & 2.7 & 4.8 \\
0.5 & 11.8 & 4.7 & 7.2 \\
0.7 & 19.7 & 7.1 & 12.6 \\
0.9 & 35.2 & 11.9 & 23.3 \\
\bottomrule
\end{tabular}

\vspace{4pt}
\begin{tabular}{c c c c}
\multicolumn{4}{l}{\textit{Panel B. Quarterly total connectedness by quantile}}\\
\toprule
$\tau$ & Overall & Contemporaneous & Lagged \\
\midrule
0.1 & 16.9 & 8.3 & 8.6 \\
0.3 & 17.6 & 6.8 & 10.8 \\
0.5 & 27.7 & 10.7 & 17.1 \\
0.7 & 55.7 & 21.6 & 34.1 \\
0.9 & 86.8 & 33.8 & 52.9 \\
\bottomrule
\end{tabular}
\begin{tablenotes}[flushleft]\footnotesize
\item \textit{Notes:} Pseudo-quantile $R^2$ (Genizi/nearPD) connectedness (\%) among the six CPI subcomponents; Overall $=$ Contemporaneous $+$ Lagged, $n_{\text{lag}}=2$. Panel A is monthly 1967--2026 (broadening ratio $\text{TCI}(0.9)/\text{TCI}(0.5)=2.98\times$); Panel B quarterly 1967Q2--2026Q2 (ratio 3.13$\times$). Both frequencies show the tail rise, increasingly driven by the lagged component. Point estimates; bootstrap confidence intervals for the six-CPI system are to be added.
\end{tablenotes}
\end{threeparttable}
\end{table}

